% GENERATED FILE. Edit canonical lessons, then run npm run generate:books.
% canonical-source-hash: fnv1a64:e93e006ca5feef73
% canonical-lessons: TA-C429-munnarivippu, TA-C429-curavali, TA-C429-alankatti, TA-C429-ataimalai, TA-C429-mekamuttam

\chapter{Storm Warnings}
\label{ch:ta-storm-warnings}
\hlchaptermodality{429}

\begin{chapteropening}
\textbf{By the end of this chapter:} \emph{I can talk about forecasts, cyclones, hail, a downpour and an overcast sky.}

Complete the last lesson of chapter 429: I can talk about forecasts, cyclones, hail, a downpour and an overcast sky.
\end{chapteropening}

\section[\texorpdfstring{\ta{முன்னறிவிப்பு} (muṉṉaṟivippu)}{முன்னறிவிப்பு (muṉṉaṟivippu)}]{\ta{முன்னறிவிப்பு} (muṉṉaṟivippu) — a forecast; advance notice}

\label{lesson:TA-C429-munnarivippu}



\begin{quote}
Before the new one: say the Tamil for to stagnate, then the Tamil for to drain away.
\end{quote}

\subsection*{You'll want to know: \ta{முன்னறிவிப்பு}}
\textbf{\ta{முன்னறிவிப்பு}} — \emph{muṉṉaṟivippu} — \textquotedblleft{}a forecast; advance notice\textquotedblright{}.

\begin{quote}
\ta{வானிலை} \ta{முன்னறிவிப்பில்} \ta{புயல்} \ta{வரும்} \ta{என்று} \ta{சொன்னார்கள்}.
\end{quote}

\begin{quote}
\emph{vāṉilai muṉṉaṟivippil puyal varum eṉṟu coṉṉārkaḷ.}
\end{quote}

\begin{quote}
In the weather forecast they said a storm would come.
\end{quote}

\ta{முன்}, before, + \ta{அறிவிப்பு}, an announcement.

\begin{cousinweb}
\textbf{\ta{முன்னறிவிப்பு}} (\emph{muṉṉaṟivippu}) is built from \ta{முன்}, \textquotedblleft{}before\textquotedblright{} (DEDR 5020), and \ta{அறிவிப்பு}, \textquotedblleft{}an announcement\textquotedblright{}, from \ta{அறிவி}, \textquotedblleft{}to make known\textquotedblright{} (DEDR 314).
\end{cousinweb}

\begin{grammarlens}[title={What you've built}]
One more word for getting around.
\end{grammarlens}

\subsection*{Your turn}
\begin{itemize}
\raggedright
  \item \emph{Say it:} \emph{muṉṉaṟivippu}
  \item \emph{Say it:} \emph{muṉṉaṟivippu}, once more
  \item \emph{Say it:} \emph{vāṉilai muṉṉaṟivippil puyal varum eṉṟu coṉṉārkaḷ.} — In the weather forecast they said a storm would come.
  \item \emph{Recall:} say the Tamil for to stagnate, then the Tamil for to drain away, then say \emph{muṉṉaṟivippu} again
\end{itemize}

\begin{tcolorbox}[breakable,colback=teal!4,colframe=teal!35!black,title={Before you move on}]
What does \ta{முன்னறிவிப்பு} mean? (\textquotedblleft{}A forecast; advance notice\textquotedblright{}.) Say it once more.
\end{tcolorbox}
\section[\texorpdfstring{\ta{சூறாவளி} (cūṟāvaḷi)}{சூறாவளி (cūṟāvaḷi)}]{\ta{சூறாவளி} (cūṟāvaḷi) — a cyclone; a whirlwind}

\label{lesson:TA-C429-curavali}



\begin{quote}
Before the new one: say the Tamil for to drain away, then the Tamil for a forecast.
\end{quote}

\subsection*{You'll want to know: \ta{சூறாவளி}}
\textbf{\ta{சூறாவளி}} — \emph{cūṟāvaḷi} — \textquotedblleft{}a cyclone; a whirlwind\textquotedblright{}.

\begin{quote}
\ta{சூறாவளியில்} \ta{பல} \ta{மரங்கள்} \ta{விழுந்தன}.
\end{quote}

\begin{quote}
\emph{cūṟāvaḷiyil pala maraṅkaḷ viḻuntaṉa.}
\end{quote}

\begin{quote}
Many trees fell in the cyclone.
\end{quote}

\ta{சூறை} alone also names a whirlwind.

\begin{cousinweb}
\textbf{\ta{சூறாவளி}} (\emph{cūṟāvaḷi}) is native Dravidian (DEDR 2743), \textquotedblleft{}a whirlwind\textquotedblright{}, beside \ta{சூறை}, which means the same; whether its end is \ta{வளி}, \textquotedblleft{}wind\textquotedblright{}, is not settled.
\end{cousinweb}

\begin{grammarlens}[title={What you've built}]
One more word for getting around.
\end{grammarlens}

\subsection*{Your turn}
\begin{itemize}
\raggedright
  \item \emph{Say it:} \emph{cūṟāvaḷi}
  \item \emph{Say it:} \emph{cūṟāvaḷi}, once more
  \item \emph{Say it:} \emph{cūṟāvaḷiyil pala maraṅkaḷ viḻuntaṉa.} — Many trees fell in the cyclone.
  \item \emph{Recall:} say the Tamil for to drain away, then the Tamil for a forecast, then say \emph{cūṟāvaḷi} again
\end{itemize}

\begin{tcolorbox}[breakable,colback=teal!4,colframe=teal!35!black,title={Before you move on}]
What does \ta{சூறாவளி} mean? (\textquotedblleft{}A cyclone; a whirlwind\textquotedblright{}.) Say it once more.
\end{tcolorbox}
\section[\texorpdfstring{\ta{ஆலங்கட்டி} (ālaṅkaṭṭi)}{ஆலங்கட்டி (ālaṅkaṭṭi)}]{\ta{ஆலங்கட்டி} (ālaṅkaṭṭi) — hailstones; hail}

\label{lesson:TA-C429-alankatti}



\begin{quote}
Before the new one: say the Tamil for a forecast, then the Tamil for a cyclone.
\end{quote}

\subsection*{You'll want to know: \ta{ஆலங்கட்டி}}
\textbf{\ta{ஆலங்கட்டி}} — \emph{ālaṅkaṭṭi} — \textquotedblleft{}hailstones; hail\textquotedblright{}.

\begin{quote}
\ta{நேற்று} \ta{மாலை} \ta{ஆலங்கட்டி} \ta{மழை} \ta{பெய்தது}.
\end{quote}

\begin{quote}
\emph{nēṟṟu mālai ālaṅkaṭṭi maḻai peytatu.}
\end{quote}

\begin{quote}
Yesterday evening there was a hailstorm.
\end{quote}

\ta{ஆலங்கட்டி} \ta{மழை} is a hailstorm.

\begin{cousinweb}
\textbf{\ta{ஆலங்கட்டி}} (\emph{ālaṅkaṭṭi}) is native Dravidian: DEDR 384 lists it, \textquotedblleft{}hail\textquotedblright{}, with \ta{ஆலம்}, \textquotedblleft{}water, rain\textquotedblright{}, and \ta{கட்டி}, \textquotedblleft{}a lump\textquotedblright{}.
\end{cousinweb}

\begin{grammarlens}[title={What you've built}]
One more word for getting around.
\end{grammarlens}

\subsection*{Your turn}
\begin{itemize}
\raggedright
  \item \emph{Say it:} \emph{ālaṅkaṭṭi}
  \item \emph{Say it:} \emph{ālaṅkaṭṭi}, once more
  \item \emph{Say it:} \emph{nēṟṟu mālai ālaṅkaṭṭi maḻai peytatu.} — Yesterday evening there was a hailstorm.
  \item \emph{Recall:} say the Tamil for a forecast, then the Tamil for a cyclone, then say \emph{ālaṅkaṭṭi} again
\end{itemize}

\begin{tcolorbox}[breakable,colback=teal!4,colframe=teal!35!black,title={Before you move on}]
What does \ta{ஆலங்கட்டி} mean? (\textquotedblleft{}Hailstones; hail\textquotedblright{}.) Say it once more.
\end{tcolorbox}
\section[\texorpdfstring{\ta{அடைமழை} (aṭaimaḻai)}{அடைமழை (aṭaimaḻai)}]{\ta{அடைமழை} (aṭaimaḻai) — a downpour; continuous rain}

\label{lesson:TA-C429-ataimalai}



\begin{quote}
Before the new one: say the Tamil for a cyclone, then the Tamil for hailstones.
\end{quote}

\subsection*{You'll want to know: \ta{அடைமழை}}
\textbf{\ta{அடைமழை}} — \emph{aṭaimaḻai} — \textquotedblleft{}a downpour; continuous rain\textquotedblright{}.

\begin{quote}
\ta{மூன்று} \ta{நாட்களாக} \ta{அடைமழை} \ta{பெய்கிறது}.
\end{quote}

\begin{quote}
\emph{mūṉṟu nāṭkaḷāka aṭaimaḻai peykiṟatu.}
\end{quote}

\begin{quote}
It has been pouring for three days.
\end{quote}

\ta{அடைமழை} is heavy rain that does not stop.

\begin{cousinweb}
\textbf{\ta{அடைமழை}} (\emph{aṭaimaḻai}) is built from \ta{அடை} and \ta{மழை}, \textquotedblleft{}rain\textquotedblright{} (DEDR 4753); how the first part should be read is not settled.
\end{cousinweb}

\begin{grammarlens}[title={What you've built}]
One more word for getting around.
\end{grammarlens}

\subsection*{Your turn}
\begin{itemize}
\raggedright
  \item \emph{Say it:} \emph{aṭaimaḻai}
  \item \emph{Say it:} \emph{aṭaimaḻai}, once more
  \item \emph{Say it:} \emph{mūṉṟu nāṭkaḷāka aṭaimaḻai peykiṟatu.} — It has been pouring for three days.
  \item \emph{Recall:} say the Tamil for a cyclone, then the Tamil for hailstones, then say \emph{aṭaimaḻai} again
\end{itemize}

\begin{tcolorbox}[breakable,colback=teal!4,colframe=teal!35!black,title={Before you move on}]
What does \ta{அடைமழை} mean? (\textquotedblleft{}A downpour; continuous rain\textquotedblright{}.) Say it once more.
\end{tcolorbox}
\section[\texorpdfstring{\ta{மேகமூட்டம்} (mēkamūṭṭam)}{மேகமூட்டம் (mēkamūṭṭam)}]{\ta{மேகமூட்டம்} (mēkamūṭṭam) — cloud cover; an overcast sky}

\label{lesson:TA-C429-mekamuttam}



\begin{quote}
Before the new one: say the Tamil for hailstones, then the Tamil for a downpour.
\end{quote}

\subsection*{You'll want to know: \ta{மேகமூட்டம்}}
\textbf{\ta{மேகமூட்டம்}} — \emph{mēkamūṭṭam} — \textquotedblleft{}cloud cover; an overcast sky\textquotedblright{}.

\begin{quote}
\ta{காலையிலிருந்து} \ta{வானம்} \ta{மேகமூட்டமாக} \ta{இருக்கிறது}.
\end{quote}

\begin{quote}
\emph{kālaiyiliruntu vāṉam mēkamūṭṭamāka irukkiṟatu.}
\end{quote}

\begin{quote}
The sky has been overcast since morning.
\end{quote}

With -\ta{ஆக} it describes the sky: \ta{மேகமூட்டமாக}, overcast.

\begin{cousinweb}
\textbf{\ta{மேகமூட்டம்}} (\emph{mēkamūṭṭam}) is built from \ta{மேகம்}, \textquotedblleft{}a cloud\textquotedblright{}, borrowed from Sanskrit \emph{megha} (CDIAL 10302), and \ta{மூட்டம்}, \textquotedblleft{}a sky overcast with clouds\textquotedblright{} (DEDR 5033).
\end{cousinweb}

\begin{grammarlens}[title={What you've built}]
One more word for getting around.
\end{grammarlens}

\subsection*{Your turn}
\begin{itemize}
\raggedright
  \item \emph{Say it:} \emph{mēkamūṭṭam}
  \item \emph{Say it:} \emph{mēkamūṭṭam}, once more
  \item \emph{Say it:} \emph{kālaiyiliruntu vāṉam mēkamūṭṭamāka irukkiṟatu.} — The sky has been overcast since morning.
  \item \emph{Recall:} say the Tamil for hailstones, then the Tamil for a downpour, then say \emph{mēkamūṭṭam} again
\end{itemize}

\begin{tcolorbox}[breakable,colback=teal!4,colframe=teal!35!black,title={Before you move on}]
What does \ta{மேகமூட்டம்} mean? (\textquotedblleft{}Cloud cover; an overcast sky\textquotedblright{}.) Say it once more.
\end{tcolorbox}
